\setcounter{paragraphcounter}{1}

\section{Beitragsordnung}

\subsection{Präambel}
Diese Beitragsordnung regelt die Mitgliedsbeiträge, Beitragsklassen und Zahlungsmodalitäten des Vereins gemäß § 14 Abs. 4 der Satzung. Sie dient der Finanzierung der Vereinsaufgaben und Aktivitäten gemäß der Vereinssatzung.

\subsection{§ \theparagraphcounter\ Erlass, Änderung, Aufhebung und Inkrafttreten} \label{bo:par:erlass_aenderung_aufhebung}
\begin{enumerate}
    \item \label{bo:abs:q650ly} Diese Beitragsordnung kann jederzeit durch die Mitgliederversammlung erarbeitet und beschlossen werden.

    \item \label{bo:abs:agu38l} Änderungen werden durch den Beschluss der Mitgliederversammlung wirksam und allen Mitgliedern bekannt gemacht.
\end{enumerate}

\refstepcounter{paragraphcounter}
\subsection{§ \theparagraphcounter\ Beitragsklassen} \label{bo:par:beitragsklassen}
\begin{enumerate}
    \item \label{bo:abs:3luqnl} Beitragsfreie Mitglieder:
    \begin{enumerate}
        \item Kinder und Jugendliche bis einschließlich des Kalenderjahres, in dem das 14. Lebensjahr vollendet wird.
        
        \item Mitglieder, die ein freiwilliges soziales Jahr (FSJ) leisten, sind auf Antrag beitragsfrei (Nachweis erforderlich).

        \item Ehrenmitglieder
    \end{enumerate}

    \item \label{bo:abs:5o2nie} Ermäßigte Beitragsklasse | 10 € pro Jahr:
    \begin{enumerate}
        \item Jugendliche ab dem 14. Lebensjahr bis zur Vollendung des 18. Lebensjahres.

        \item Jugendliche unter 14 Jahren, die bereits ein vollwertiges Häs besitzen.
        
        \item Schüler und Studenten (Nachweis erforderlich).
    \end{enumerate}

    \item \label{bo:abs:ofkmx8} Vollzahler | 30 € pro Jahr:
    \begin{enumerate}
        \item Mitglieder ab Vollendung des 18. Lebensjahres.
        
        \item Auszubildende oder Personen in Fortbildung mit Entlohnung.
    \end{enumerate}
\end{enumerate}

\refstepcounter{paragraphcounter}
\subsection{§ \theparagraphcounter\ Zahlungsmodalitäten} \label{bo:par:zahlungsmodalitaeten}
\begin{enumerate}
    \item \label{bo:abs:0dsvko} Der Beitrag und das Sprunggeld werden im SEPA-Lastschriftverfahren eingezogen. Jedes Mitglied ist verpflichtet, dem Verein hierzu eine gültige Einzugsermächtigung zu erteilen.

    \item \label{bo:abs:j3l19p} Fehlschlagende Lastschriften
    \begin{enumerate}
        \item Schlägt eine Lastschrift aus Gründen fehl, die das Mitglied zu vertreten hat (z.B. mangelnde Kontodeckung, falsche Bankdaten), trägt das Mitglied die dem Verein dadurch entstehenden Bank- und Verwaltungskosten.

        \item Zusätzlich erhebt der Verein eine Bearbeitungsgebühr von 5,00 € pro fehlgeschlagener Lastschrift.

        \item Der ausstehende Beitrag ist anschließend innerhalb von 14 Tagen nach schriftlicher Mitteilung des Vereins auf das Vereinskonto zu überweisen.

        \item Kommt ein Mitglied trotz Mahnung seinen Zahlungsverpflichtungen nicht nach, gelten die Regelungen der Satzung.

        \item Auf Antrag kann der Kassenwart in begründeten Fällen Ratenzahlung oder Stundung des Beitrags gewähren.
    \end{enumerate}
\end{enumerate}

\refstepcounter{paragraphcounter}
\subsection{§ \theparagraphcounter\ Beitragsänderung} \label{bo:par:beitragsaenderung}
\begin{enumerate}
    \item \label{bo:abs:6lv8zl} Erhöht/Senkt die Mitgliederversammlung die Beiträge, ist die Erhöhung/Senkung ab Beginn des auf die Beschlussfassung folgenden Kalenderjahres wirksam, sofern die Mitgliederversammlung nichts anderes beschließt.

    \item \label{bo:abs:xgnm78} Eine Rückwirkung auf bereits gezahlte Beiträge ist ausgeschlossen.
\end{enumerate}

\refstepcounter{paragraphcounter}
\subsection{§ \theparagraphcounter\ Inkrafttreten} \label{bo:par:inkrafttreten}
\begin{enumerate}
    \item \label{bo:abs:9v9j5x} Diese Beitragsordnung tritt nach Beschluss der Mitgliederversammlung vom dd. MMMM YYYY in Kraft. \\
    Sie ersetzt alle vorherigen Versionen.
\end{enumerate}